\documentclass[10pt,a4paper]{amsart}
\usepackage[margin=19mm]{geometry}
\usepackage{amsmath,amssymb,mathtools}
\usepackage[scr=rsfs,cal=euler]{mathalfa}
\usepackage{xfrac}
\usepackage[unicode,hidelinks,bookmarks=false]{hyperref}
\newcommand{\ie}{i.e.\ }
\newcommand{\cf}{cf.\ }
\newcommand{\resp}{respectively\ }
\newcommand{\Subsection}{\subsection}
\providecommand{\nobreakdash}{\nobreak\hbox{-}}
\setlength{\emergencystretch}{2em}
\multlinegap0pt
\usepackage{amssymb}
\usepackage{tikz}
\usepackage[shortlabels]{enumitem}
\usepackage[normalem]{ulem}
\newtheorem{theorem}{Theorem}
\newtheorem{lemma}{Lemma}
\def\J{\mathcal J}
\title{More minor summation formulae}
\author{Shane Chern, Theresia Eisenk\"olbl, Ilse Fischer, Moritz Gangl, Mona Gatzweiler, \'Alvaro Guti\'errez, Christian Krattenthaler, Nishu Kumari, Markus Reibnegger, Marcus Sch\"onfelder, Atsuro Yoshida}
\date{Source: arXiv:2603.21021v1 (CC BY 4.0)}
\begin{document}
\maketitle

\noindent\emph{Source and licence.} More minor summation formulae. Version arXiv:2603.21021v1 (CC BY 4.0), distributed by arXiv under the Creative Commons Attribution 4.0 International licence, http://creativecommons.org/licenses/by/4.0/, as recorded in the arXiv licence metadata for this exact version and shown on the abstract page https://arxiv.org/abs/2603.21021v1. Full licence notice: \texttt{licenses/arxiv-2603.21021-CC-BY-4.0.txt}.\par\smallskip

\noindent\textit{Editorial scope.} Counted unit: the article's lemma lem:X2 (source0.tex lines 1132--1151). The article's complete original proof of it is delivered with it (source0.tex lines 1153--1223, one \texttt{proof} environment of 71 lines). No mathematics is written, repaired or re-derived: the statement and the proof are the article's own lines, in the article's own notation, and no retained line is substituted or re-flowed.  Retained uncounted as the source-local context the proof needs: lines 283--316; lines 1066--1085.  Nothing was omitted from any retained span, so the package carries no omission record.  No retained line contains a citation, so the wrapper carries no bibliography and no bibliography entry.  No retained line contains a figure environment, an image-inclusion command or drawing code, so no image is delivered and the package declares no asset dependency.  Editorial formatting only, in addition: an AMS \texttt{amsart} preamble that restates the article's own declarations and macros used by the retained lines, namely the environments \texttt{theorem}, \texttt{lemma} on separate counters, together with the standard packages those lines need.  The retained environments print in this document as Theorem 1, Lemma 1 in order of appearance, whereas the article prints the same items with the numbering of its own full document.  The complete native source of the article as distributed in the arXiv e-print for this exact version is bundled unchanged as \texttt{sources/196910/source0.tex}.
\begin{theorem} \label{thm:main1}
Let $A$ and $B$ be $m\times n$ matrices,
$X$ an $n\times n$ matrix, and
$\J =(1)_{1\le i,j\le n}$. If $m$ is even, we have
\begin{equation} \label{eq:1}
\det\bigl(AXB^t+B(\J -X^t)A^t\bigr)
=f_{A,B}(X)f_{B,A}(\J -X^t)
\end{equation}
where
$$
f_{A,B}(X)=\underset {|I|=|J|=m/2}{\sum_{I,J\subseteq[n]}}
\det X_{I,J} \det(A^IB^J),
$$
with $X_{I,J}$ denoting the submatrix of~$X$ consisting of the entries
in the 
rows indexed by~$I$ and columns indexed by~$J$, $A^I$ denoting
the submatrix of~$A$ consisting of the columns indexed by~$I$,
with an analogous meaning of~$B^J$, and with $A^IB^J$ denoting
the concatenation of the matrices $A^I$ and $B^J$.

If $m$ is odd, we have
\begin{equation} \label{eq:1o}
\det\bigl(AXB^t+B(\J -X^t)A^t\bigr)
=g_{A,B}(X)g_{B,A}(\J -X^t)
=(-1)^{(m-1)/2}g_{A,B}(X)g_{B,A}(X^t),
\end{equation}
where
$$
g_{A,B}(X)=\underset {|I|=(m+1)/2,\ |J|=(m-1)/2}{\sum_{I,J\subseteq[n]}}
\det\bigl(\mathbf1 \, X_{I,J}\bigr) \det(A^IB^J),
$$
where $\mathbf 1$ denotes a column of length $(m+1)/2$ consisting
entirely of~$1$'s.
\end{theorem}

\begin{lemma} \label{lem:X1}
Let $n$ and $\ell$ be positive integers.
Let $X$ be an upper triangular $n\times n$ matrix with $1$'s above
the diagonal, that is, $X_{i,j}=1$ for $1\le i<j\le n$.
Furthermore let $I=\{i_1,i_2,\dots,i_\ell\}$ and~$J=\{j_1,j_2,\dots,j_\ell\}$
be subsets of~$[n]$ of equal cardinality, both of which we assume to
be ordered.
Then, with the notation of Theorem~\ref{thm:main1}, we have
\begin{equation} \label{eq:X1} 
\det(X_{I,J})=\begin{cases}
\prod _{k=1} ^{\ell}X_{i_k,i_k}^{\chi(i_k=j_k)}
(1-X_{i_k,i_k})^{\chi(j_{k-1}=i_k)},&
\text{if }1\le i_1\le j_1\le i_2\le\dots\le j_\ell\le n,\\
0,&\text{otherwise}.
\end{cases}
\end{equation}
Here, by convention, $j_0=0$, and
$\chi(\mathcal S)=1$ if $\mathcal S$ is true and
$\chi(\mathcal S)=0$ otherwise.
\end{lemma} 

\begin{lemma} \label{lem:X2}
Let $n$ and $\ell$ be positive integers.
Let $X$ be an upper triangular $n\times n$ matrix with $1$'s above
the diagonal, that is, $X_{i,j}=1$ for $1\le i<j\le n$.
Furthermore let $I=\{i_1,i_2,\dots,i_{\ell+1}\}$
and~$J=\{j_1,j_2,\dots,j_\ell\}$
be subsets of\/~$[n]$, both of which we assume to
be ordered.
Then, with the notation of Theorem~\ref{thm:main1}, we have
\begin{equation} \label{eq:X2} 
\det(\mathbf 1\,X_{I,J})=\begin{cases}
(-1)^\ell\prod _{k=1} ^{\ell}X_{i_k,i_k}^{\chi(i_k=j_k)}
(1-X_{i_{k+1},i_{k+1}})^{\chi(j_{k}=i_{k+1})},&\\
&\kern-2cm
\text{if }1\le i_1\le j_1\le i_2\le\dots\le j_\ell\le i_{\ell+1}\le n,\\
0,&\kern-2cm\text{otherwise}.
\end{cases}
\end{equation}
Here, $\mathbf1$ denotes a column of $\ell+1$ elements~$1$.
\end{lemma} 

\begin{proof}
We prove the claimed identity by induction on~$\ell$, also using
the result from Lemma~\ref{lem:X1}.

The claim is trivially true for $\ell=0$ (if we interpret the empty
product as~1)
and for $\ell=1$. Therefore, from now on, we assume that $\ell\ge2$.

For the induction step, we distinguish several cases.

\begin{itemize}
\item If $i_{\ell+1}= j_\ell$, then the last row of~$X_{I,J}$ 
contains only one non-zero element, namely
$X_{i_{\ell+1},j_\ell}=X_{i_{\ell+1},i_{\ell+1}}$
as the last entry in the row. By expansion along the last row,
we get
\begin{equation} \label{eq:8} 
\det(\mathbf1\,X_{I,J}) =
(-1)^{\ell}\det(X_{I\backslash\{i_{\ell+1}\},J})+
X_{ i_{\ell+1} ,i_{\ell+1}}
\det(\mathbf1\, X_{ I\backslash
\{ i_{\ell+1} \}, J\backslash \{ j_\ell \} }).
\end{equation}
Since $i_{\ell} <  i_{\ell+1}=j_\ell$, by Lemma~\ref{lem:X1} and
by induction hypothesis, we obtain
\begin{align*}
\det(\mathbf1\,X_{I,J}) &=(-1)^{\ell}
\prod _{k=1} ^{\ell}X_{i_k,i_k}^{\chi(i_k=j_k)}
(1-X_{i_k,i_k})^{\chi(j_{k-1}=i_k)}\\
&\kern1cm
+
X_{ i_{\ell+1} ,i_{\ell+1}}
(-1)^{\ell-1}\prod _{k=1} ^{\ell-1}X_{i_k,i_k}^{\chi(i_k=j_k)}
(1-X_{i_{k+1},i_{k+1}})^{\chi(j_{k}=i_{k+1})}\\
&=(-1)^{\ell}(1-X_{i_{\ell+1},i_{\ell+1}})
\prod _{k=1} ^{\ell}X_{i_k,i_k}^{\chi(i_k=j_k)}
(1-X_{i_{k},i_{k}})^{\chi(j_{k-1}=i_{k})}
\end{align*}
for the case where $i_1\le j_1\le j_2\le\dots\le j_\ell$.
If that condition is violated,
we get zero for both summands on the right-hand side of~\eqref{eq:8},
again due to Lemma~\ref{lem:X1} and the induction hypothesis.
This is in accordance with~\eqref{eq:X2}.
\item
If $i_{\ell+1} < j_\ell$, then the last 
column of~$X_{I,J}$ is a column of~1's. Consequently, the determinant
$\det(\mathbf1\,X_{I,J})$ vanishes, as
desired.
\item
If $i_{\ell+1} > j_\ell$ then the last row of $X_{I,J}$ is a row of 0's.
We may therefore expand the determinant $\det(\mathbf1\,X_{I,J})$
along the last row and obtain
$$
\det(\mathbf1\,X_{I,J})=(-1)^{\ell}
\det(X_{I\backslash\{i_{\ell+1}\},J}).
$$
If $i_1\le j_1\le i_2\le \dots\le j_\ell$, then, by
Lemma~\ref{lem:X1}, we infer
$$
\det(\mathbf1\,X_{I,J})=(-1)^{\ell}
\prod _{k=1} ^{\ell}X_{i_k,i_k}^{\chi(i_k=j_k)}
(1-X_{i_{k},i_{k}})^{\chi(j_{k-1}=i_{k})},
$$
which is in accordance with \eqref{eq:X2} since $j_\ell<i_{\ell+1}$.
If the above inequality chain does not hold then, again by
Lemma~\ref{lem:X1}, the determinant
$\det(X_{I\backslash\{i_{\ell+1}\},J})$ vanishes, and hence also
$\det(\mathbf1\,X_{I,J})$.
\end{itemize}
This completes the proof of the lemma.
\end{proof}

\end{document}
